%%%%%%%%%%  scan idx 85  |  folio 72  %%%%%%%%%%
variété \(Y\) de \(X\)), on note \(\gamma_X(\underline F)\) (resp.\ \(\gamma_X(E)\), resp.\ \(\gamma_X(Y)\)) l'élément de \(K(X)\) correspondant à \(\underline F\) (resp.\ à \(\underline E\), resp.\ à \(\mathcal O_Y\)). (En fait, on écrit souvent \(\underline F\), \(E\) au lieu de \(\gamma_X(\underline F)\)\dots).

On note \(A(X)\) l'anneau gradué des classes de cylces de \(X\) (pour l'équivalence rationnelle -cf.\ Chow)~; \(A^i(X)\) est formé des cycles de codimension \(i\).

Si \(E\) est un fibré à fibre vectoriel, on note \(C(E)=\sum C^i(E)\), \(C^i(E)\in A^i(X)\), sa classe de Chern (définie comme image réciproque des cycles de Schubert)~; on en déduit par linéarité les \(C^i(y)\), pour tout \(y\in K(X)\). Si \(C(y)=\prod(1+a_i)\), on pose suivant Hirzebruch~:
\[
  \ch(y)=\sum e^{a_i},\qquad
  T(y)=\prod\frac{a_i}{1-e^{-a_i}}.
\]
Ce sont des éléments de \(A(X)\otimes\mathbb Q\).

Lorsque \(y\) est le fibré tangent à \(X\), on écrit \(T(X)\) au lieu de \(T(y)\).

Si \(G\) est un fibré à fibre vectorielle, on note \(\lambda^pG\) sa puissance extérieure \(p\)-ème, et on pose \(\lambda_{-1}(G)=\sum(-1)^p\lambda^p(G)\)~; c'est un élément de \(K(X)\).

Si \(f\) est une application régulière propre d'une variété \(Y\) dans une variété \(X\), on définit des homomorphismes (additifs)~:
\[
  f^*:A(X)\longrightarrow A(Y)\qquad\text{et}\qquad
  f_*:A(Y)\longrightarrow A(X).
\]

Le premier est homogène de degré \(0\) et multiplicatif. Le second est homogène de degré \(\dim.X-\dim.Y\).

On définit également des homomorphismes~:
\[
  f^!:K(X)\longrightarrow K(Y)\qquad\text{et}\qquad
  f_!:K(Y)\longrightarrow K(X).
\]

Si \(\underline F\) est un faisceau cohérent sur \(X\), \(f^!(\underline F)\) est égal, par définition, à la somme alternée des \(\Tor_p^{\mathcal O_X}(\mathcal O_Y,\underline F)\). Pour les fibrés à fibre vectorielle, c'est simplement l'opération d'image réciproque.

